\section{\hspace{-0.5em}SU Cross-Layer Optimization for Coded Transmission}
\label{sec:3}

Here, we present a \myCLO{} approach for the \mySU{} \myLDPC{} coded transmission equipped with \myJBPL{} and \myCA{}. The optimization is considered using \myCC{} subject to \myQoS{} constraints in the presence of a PU transmission, thermal noise, and \myTPI{}. In Sect.~\ref{sec:3A}, we derive an upper bound for the coded \mySU{} bit-error probability (\myBEP{}). This bound is used to define the \mySU{} packet-error probability (\myPEP{}). In Sect.~\ref{sec:3c}, we employ a \myTHARQ{} to maximize a lower bound on the \mySU{} \mySE{}. Subsequently, the time-slot index $\ell$ is dropped as the channel is considered time-invariant.

\subsection{Physical Layer Parameters}
\label{sec:3A}

In~(\ref{eq:1}), we define the random signal-to-noise-plus-interference ratio (\mySINR{}) of the $\myk$th subcarrier of the \mySU{} received signal at the output of the combiner after $\myJ=\myj$ transmissions as
\begin{equation}
\label{eq:2}
\myA^{(\myj)}_{\mySrm,\myk} =
\myp_\myk \, |\myD_\myk|^2 \;
\frac{
    \Bigl[ \sum_{{\myj}'=1}^{\myj} |\myHtilde_{\mySrm\mySrm,\myk,{\myj}'}|^2 \Bigr]^2
}{
    \sum_{{\myj}'=1}^{\myj} |\myHtilde_{\mySrm\mySrm,\myk,{\myj}'}|^2 
    \Bigl( 1 + \myp_{\myPrm} \, |\myHtilde_{\myPrm\mySrm,\myk,{\myj}'}|^2 \Bigr)
}
\end{equation}

Here, $\myp_{\myPrm}$ is the average \myPU{} transmit power per subcarrier and $|\myD_\myk|$ is the magnitude of the \myQAM{} symbol. Also, $\myHtilde_{\mySrm\mySrm,\myk,\myj}$ and $\myHtilde_{\myPrm\mySrm,\myk,\myj}$ represent the normalized \myTF{} channel coefficients between the \mySU{} \myTX{} and \mySU{} \myRX{}, and between the \myPU{} \myTX{} and \mySU{} \myRX{}, respectively, for the $\myk$th subcarrier, where

\begin{align}
\tilde{H}_{\mySrm\mySrm,\myk,\myj} &\coloneqq \frac{H_{\mySrm\mySrm,\myk,\myj}}{\sigma_{\mySrm,\myk}}, 
& \text{where} \;\; \tilde{H}_{\mySrm\mySrm,\myk,\myj} \sim \mathcal{CN}(0,1), \\[0.3em]
\tilde{H}_{\myPrm\mySrm,\myk,\myj} &\coloneqq \frac{H_{\myPrm\mySrm,\myk,\myj}}{\sigma_{\mySrm,\myk}}, 
& \text{where} \;\; \tilde{H}_{\myPrm\mySrm,\myk,\myj} \sim \mathcal{CN}(0,1).
\end{align}

with $\myCcal\myNcal(\mu;\sigma^{2})$ representing a circularly-symmetric complex Gaussian random variable with mean $\mu$ and variance $\sigma^{2}$. 
It is noted that $\bigl|\myHtilde_{\mySrm\mySrm,\myk,\myj}\bigr|^{2}$ and $\bigl|\myHtilde_{\myPrm\mySrm,\myk,\myj}\bigr|^{2}$ follow central chi-square distributions as normalized channel gains, respectively:

\begin{align}
\mathrm{Var}\bigl[\tilde{H}_{\mySrm\mySrm,\myk,\myj}\bigr] 
&= \mathbb{E}\Bigl[ \bigl|\tilde{H}_{\mySrm\mySrm,\myk,\myj}\bigr|^2 \Bigr] 
= \frac{\sigma^2_{\mySrm\mySrm}}{\sigma^2_{\mySrm,\myk}}, \\[0.3em]
\mathrm{Var}\bigl[\tilde{H}_{\myPrm\mySrm,\myk,\myj}\bigr] 
&= \mathbb{E}\Bigl[ \bigl|\tilde{H}_{\myPrm\mySrm,\myk,\myj}\bigr|^2 \Bigr] 
= \frac{\sigma^2_{\myPrm\mySrm}}{\sigma^2_{\mySrm,\myk}}.
\end{align}


Here, $\sigma^{2}_{\mySrm,\myk} = \sigma^{2}_{\mySrm,\myImath,\myk} + \sigma^{2}_{\mySrm,\myW}$ represents the average power of \myTPI{} plus noise at the $\myk$h subcarrier of the \mySU\ \myRX{}.\\ Using the channel gains and after some algebra, the probability density function (\myPDF{}) $\myp_{\myA^{(\myj)}_{\mySrm,\myk}} (\alpha)$ of $A^{(\myj)}_{\mySrm,k}$ in (\ref{eq:4}) conditioned on $\left|\myD_\myk\right| = \myd_\myk$ is defined as
\begin{align}
f^{(\myj)}_{\myA_{\mySrm,\myk}}(\alpha) 
&:= \myu(\alpha) \; 
   \frac{\Gamma(3\myj)}{\Gamma(\myj)^2} 
   \left( \frac{\kappa_\myk}{8} \right)^{\myj/2} 
   \alpha^{\myj/2 - 1} 
   \exp\Bigl(\frac{\kappa_\myk}{8} \, \alpha \Bigr) 
   \nonumber \\[0.2em]
&\quad \times \mathcal{D}_{-3\myj}\Bigl( 2 \sqrt{\frac{\kappa_\myk}{8} \, \alpha} \Bigr),
\label{eq:3}
\end{align}

\noindent with 
\[
\kappa_\myk = \frac{1}{\myp_\myk \, \myd_\myk \, \sigma^2_{\mySrm\mySrm,\myk}}, \quad
\mathcal{D}_\nu(\cdot) \text{ denotes the parabolic cylindrical function~\cite{Gradshteyn2007_alt}}.
\]

and $u(\cdot)$ represents the unit-step function.

To evaluate the system performance, we need the coded \mySU\ transmission \myBEP{}, a closed form representation of which is not available. Instead, we use an upper bound on the expected coded \myBEP{} $\myP^{(\mycorm)}_{\mybrm,\mathrm{\myQAM{}},\myj}$ of the $\myj$th transmission for 
$\myA^{(\myj)}_{\mySrm,\myk} = \alpha^{(\myj)}_{\mySrm,\myk}$ averaged over the statistics of 
$\mathbf{\myA}^{(\myj)}_{\mySrm} = \big[\, \alpha^{(\myj)}_{\mySrm,0},\ldots,\alpha^{(\myj)}_{\mySrm,\myK-1} \,\big]^{\myTrm}$.

Averaging over $\myK$ subcarriers, one obtains

\begin{align}
\myP^{(\mycorm)}_{\mybrm,\myQAMrm,\myj}
&\le \myP^{(\myj)}_{\myUBrm} \notag \\
&\hspace{-1cm}= \sum_{\myd = \myd_{\min}}^{\myn}
   \sum_{\myw = 1}^{\myk}
   \frac{\myw}{\myk}\, \myA_{\myw,\myd}
   \sum_{\nu = 0}^{\lfloor\frac{\myd}{2}\rfloor}
   \binom{\lfloor\frac{\myd}{2}\rfloor}{\nu}
   (-1)^{\nu}\,
   \xi(\nu,\myd,\myj),
\label{eq:4}
\end{align}

with \!$\xi(\nu,\myd,\myj)
=\!\myE[\myAbf{}^{(\myj)}_{\mySrm}]
{
\Bigl(\myP^{(\myunrm)}_{\mybrm,\myQAMrm},\myj\Bigr)^{\lfloor \tfrac{\myd}{2} \rfloor + \nu}
}.$ Here, $\myd_{\myminmathrm}$ and $\myA_{\myw,\myd}$ denote the minimum distance of the parity check matrix corresponding to an $(\myn,\myk)$ \myLDPC\ code, and the number of possible code words of Hamming weight $\myd$ with $\myw$ information bits\linebreak in error, respectively. Furthermore, $\myP^{(\myunrm)}_{\mybrm,\myQAMrm,\myj}$ is the conditional uncoded \mySU\ \myBEP{} of the $\myj$th transmission averaged over $\myK$ subcarriers \cite{Mutti2003_alt}, and is given by
\begin{equation*}
P^{(\myunrm)}_{\mybrm{},\myQAMrm,\myj}
= \frac{1}{3\myMtot}
  \sum_{\myk \in \Omega_{0,\myK}} \beta_{\myk}\,
  \myf\bigl(\alpha^{(\myj)}_{\mySrm,\myk}\bigr)
= \frac{\myG\bigl(\myAbf{}^{(\myj)}_{\mySrm}\bigr)}{3\myMtot}.
\end{equation*}
Here, $\beta_{\myk}= \frac{\sqrt{\myM_{\myk}} - 1}{\sqrt{\myM_{\myk}}}$, with
$\myM_{\myk} = 2^{\mym_{\myk}}$ and the function $
\myf\bigl(\alpha^{(\myj)}_{\mySrm,\myk}\bigr)
=
e^{-\frac{3\alpha^{(\myj)}_{\mySrm,\myk}}{2(\myM_{\myk}-1)}}
+ 3 e^{-\frac{2\alpha^{(\myj)}_{\mySrm,\myk}}{\myM_{\myk}-1}}
$ is the approximation of the Q-function as a sum of two weighted exponential functions~\cite{Mutti2003_alt,Karagiannidis2007_alt}.
The sum of $\myK$ i.i.d. random variables can be expressed as
\[
\mathcal{G}\Bigl(\mathbf{A}^{(\myj)}_{\mySrm}\Bigr) 
= \sum_{k \in \Omega_{0,K}} \beta_k \; 
  f\Bigl(A^{(\myj)}_{\mySrm,k}\Bigr).
\]
\newpage

\begin{minipage}[t][12cm]{\columnwidth}

is approximated as a Gaussian distributed RV
\begin{equation*}
\mathcal{G}\Bigl(\mathbf{A}^{(\myj)}_{\mySrm}\Bigr) \sim 
\mathcal{N}\Bigl( \mu_{\mathcal{G},\myj}, \; \sigma^2_{\mathcal{G},\myj} \Bigr),
\end{equation*}

\noindent
where
\begin{equation*}
\mu_{\mathcal{G},\myj} = \sum_{k \in \Omega_{0,K}} \mu_{k,\myj}, 
\quad \text{and} \quad
\sigma^2_{\mathcal{G},\myj} = \sum_{k \in \Omega_{0,K}} \sigma^2_{k,\myj}.
\end{equation*}

Here, the mean and variance for subcarrier $k$, i.e., $\mu_{k,\myj}$ and $\sigma^2_{k,\myj}$, 
are computed by averaging 
$\beta_k f\bigl(A^{(\myj)}_{\mySrm,k}\bigr)$ 
over the joint statistics of $A^{(\myj)}_{\mySrm,k}$ and the QAM symbol power $|D_k|^2$, 
where $|D_k|^2 = \epsilon$ has the probability mass function~\cite{Mansour2018_alt}:

\begin{equation}
p_{|D_k|^2}(\epsilon)
=
\sum_{i=1}^{K_{m_k}} S_{i,m_k} \;
\delta\Bigl(\epsilon - d^2_{i,m_k} \Bigr),
\label{eq:5}
\end{equation}

with 
$K_{m_k} = 3 - 2 m_k + \frac{m_k^2}{2}$ 
distinct power levels $|D_k|^2 = d^2_{i,m_k}$ in a single quadrant of $\Omega_{\myQAMrm}$.  
The sets $\{ d^2_{i,m_k} \} \subset \Omega_{d,m_k}$ 
and $\{ S_{i,m_k} \} \subset \Omega_{S,m_k}$ are given in~\cite{Ding2012_alt}.  

Using (\ref{eq:3}) and (\ref{eq:5}), the mean can be written as

\begin{equation*}
\begin{split}
\mu_{k,\myj} 
&= \mathbb{E}\bigl[ \beta_k f(A^{(\myj)}_{\mySrm,k}) \bigr] 
= \beta_k \int_0^{\infty} f(\alpha) \, p_{A^{(\myj)}_{\mySrm,k}}(\alpha) \, d\alpha \\[0.2em]
&= \beta_k \frac{\Gamma(3\myj)}{\Gamma(\myj)^2} 
  \sum_{i=1}^{K_{m_k}} S_{i,m_k} 
  \Bigl( \frac{\kappa_{k,i}}{8} \Bigr)^{\myj/2} 
  \int_0^{\infty} 
    \alpha^{\myj/2-1} \,
    \mathcal{D}_{-3\myj}\Bigl( 2 \sqrt{\frac{\kappa_{k,i} \alpha}{8}} \Bigr) 
    e^{\kappa_{k,i}\alpha/8} \\ 
&\quad \times \Bigl( e^{-c_{1,k} \alpha} + 3 \, e^{-c_{2,k} \alpha} \Bigr) 
\, d\alpha \\[0.2em]
&= \beta_k \, \theta_{\myj} 
  \sum_{i=1}^{K_{m_k}} S_{i,m_k} 
  \Bigl( \frac{\kappa_{k,i}}{8} \Bigr)^{\myj/2} 
  \Bigl( \Psi(c_{1,k},\myj) + 3 \, \Psi(c_{2,k},\myj) \Bigr).
\end{split}
\end{equation*}

Here, the function $\Psi(x,\myj)$ is defined as
\[
\Psi(x,\myj) = x^{-\myj/2} \;
\mathcal{F}\Bigl( \tfrac{3\myj}{2}, \tfrac{\myj}{2}; 2\myj + \tfrac{1}{2}; 1 - \tfrac{\kappa_{k,i}}{4x} \Bigr),
\]
where $\mathcal{F}(a,b;c;z)$ denotes the hypergeometric function~\cite{Gradshteyn2007_alt}, 
and 
\[
\kappa_{k,i} = \frac{1}{p_k \, d^2_{i,m_k} \, \sigma^2_{\mySrm\mySrm,k}}.
\]

The constants are given by
\[
c_{1,k} = \frac{3}{2(M_k-1)}, \quad
c_{2,k} = \frac{2}{M_k-1}, \quad
\theta_{\myj} = 
\frac{\Gamma(3\myj) \, \sqrt{\pi}}{\Gamma(\myj) \, \Gamma\bigl(2\myj + \tfrac{1}{2}\bigr) \, (\sqrt{2})^{5\myj-2}}.
\]

Similarly, the variance for subcarrier $k$ is obtained as
\[
\sigma^2_{k,\myj} = 
\mathbb{E}\Bigl[ \bigl( \beta_k f(A^{(\myj)}_{\mySrm,k}) \bigr)^2 \Bigr] - \mu^2_{k,\myj}.
\]
can be found using

\begin{equation*}
\hspace*{-1cm}
\myE[\myA^{(\myj)}_{\mySrm,\myk}]
{
   \left(\myf\left(\myA^{(\myj)}_{\mySrm,\myk}\right)\right)^2
}
= \theta_\myj \sum_{\myi=1}^{\myK_{\mym_\myk}}
\myS_{\myi,\mym_\myk}
\left(\frac{\kappa_{\myk,\myi}}{8}\right)^{\frac{\myj}{2}}
\end{equation*}

\begin{equation*}
\hspace*{0.5cm}
\times\quad \left(
   \Psi(\myc_{3,\myk}, \myj)
   + 9\Psi(\myc_{4,\myk}, \myj)
   + 6\Psi(\myc_{5,\myk}, \myj)
\right),
\end{equation*}

where
$
\myc_{3,\myk} = \frac{3}{\myM_\myk - 1},
\
\myc_{4,\myk} = \frac{4}{\myM_\myk - 1},
\myc_{5,\myk} = \frac{7}{2(\myM_\myk - 1)}.
$

Finally, the upper bound in (\ref{eq:4}) is given by

\begin{equation*}
\begin{aligned}
\myP^{(\myj)}_{\myUBrm}
&=
\sum_{\myd=\myd_{\min}}^{\myn}
\sum_{\myw=1}^{\myk}
\frac{\myw}{\myk}\, \myA_{\myw,\myd}
\sum_{\nu=0}^{\lfloor \frac{\myd}{2} \rfloor}
\binom{\lfloor \frac{\myd}{2} \rfloor}{\nu}
(-1)^{\nu}
\\[4pt]
&\hspace{0.5cm}\times
\left( \frac{1}{3\myMtot} \right)^{\myp}
\sum_{\myi=0}^{\myp}
\binom{\myp}{\myi}\,
\mu^{\myi}_{\myG,\myj}\,
\zeta_{\myp-\myi,\myj},
\end{aligned}
\end{equation*}

\end{minipage}

\newpage
\begin{minipage}[t][12cm]{\columnwidth}

with $\myp = \lfloor \frac{\myd}{2} \rfloor +\nu$ and the central moments for $\myp \in \myNbb$ given by
\begin{align*}
\zeta_{\myp,\myj}
&=
\quad\myE[\myAbf^{(\myj)}_{\mySrm}]
{
    \left(
        \myG\left( \myAbf^{(\myj)}_{\mySrm} \right)
        - \mu_{\myG,\myj}
    \right)^{\myp}
}
\\
&=\quad
1 \cdot 3 \cdot \ldots (\myp-1)\,
\sigma_{\myG,\myj}^{\,\myp}\,
( (\myp+1)\bmod\!2 ).
\end{align*}
The dependeny of the different expressions on $\myR_\myc,$ $\mypbf$ and $\mymbf$ \!is not indicated for better readability of the resulting formulas.
\vspace{-1em}
\subsection{Data Link Layer Parameters}
\label{sec:3b}

The \mySU{} \myRX{} determines if at least one estimated bit $\hat \myB_{\kappa}$ in the received packet is in error compared to the information bit
$B_{\kappa}$, and generates a \myTHARQ{} over a one-bit feedback channel
to request the \mySU{} \myTX{} to retransmit the packet. The \mySU{} \myTX{}
\myTHARQ{} determines and retransmits the required packet from
the \myTX{} buffer. The \myPEP{} $\myP^{(\myj)}_{\mypac}$ of the \mySU{} transmission after $\myj$ transmissions can be upper bounded by a union bound $\myP^{(\myj)}$ given by
\[
\myP^{(\myj)}_{\mypacrm}
= 
\Pr\left\{
    \bigcup_{\kappa = 1}^{\myL \myMtot\myR_{\myc}}
    \; \myB_{\kappa} \ne \hat{\myB}_{\kappa}
\right\}
\le
\myL \myMtot \myR_{\myc} \, \myP^{(\myj)}_{\myUBrm}
= 
\myP^{(\myj)} ,
\]
with 
$\myP^{(\myj)} = \myP^{(\myj)}(\mypbf, \mymbf, \myR_{\myc})$.
To simplify the system analysis,
the following assumptions have been made:
\begin{itemize}[leftmargin=1.9em]
    \renewcommand\labelitemi{$\scriptstyle\bullet$}
    \item A packet is deemed lost, if it is still decoded with errors after $\myJ_{\max}+1$ transmissions.
    \item The noise and \myTPI\ profiles are known at the \mySU\ \myRX\ and \myPU\ \myRX{}.
    \item Perfect error detection is achieved using \myCRC{}.
    \item The feedback channel is error-free and has no latency.
\end{itemize}
To evaluate the system performance of the \mySU\ transmission,
the average \mySU\ \mySE\ is defined by
\begin{equation}
\bar{\myS}(\mypbf,\mymbf,\myR_\myc)
  = \frac{\myL \myMtot \myR_\myc}{\myK}
    \, \myE{\frac{\myImath(\myJ)}{\myJ}},
\label{eq:6}
\end{equation}
with the indicator function
\[
\myImath(\myj) =
\begin{cases}
  1, & \text{for\quad} 1 \le \myj \le \myJ_{\max}+1 \\[0.3em]
  0, & \text{else}.
\end{cases}
\]
Furthermore, the probability for a packet to require $\myj$ transmissions for error-free reception is given by
\begin{equation}
\Pr\{\myj\} =
\begin{cases}
\displaystyle \Bigl( 1 - \myP_{\mypacrm}^{(\myj)} \Bigr) 
\prod_{i=1}^{\myj-1} \myP_{\mypacrm}^{(i)}, 
& 1 \le \myj \le \myJ_{\max}, \\[1em]
\displaystyle \prod_{i=1}^{\myj-1} \myP_{\mypacrm}^{(i)}, 
& \myj = \myJ_{\max}+1.
\end{cases}
\label{eq:7}
\end{equation}

\noindent
Using (\ref{eq:6}) and (\ref{eq:7}), a lower bound on the spectral efficiency can be written as
\[
\myS_{\myLBrm} \le \bar{\myS}(\mypbf,\mymbf,\myR_\myc),
\]
where
\begin{equation}
\myS_{\myLBrm} 
= \alpha \, \mathbb{E} \Biggl[
    \frac{\myI(\myJ)}{\myJ} \;\Big|\; \myP_{\mypacrm}^{(\myj)} \to \myP^{(\myj)}
\Biggr]
= \alpha \left(
    1 - \sum_{\myj=1}^{\myJ_{\max}} 
        \frac{\displaystyle \prod_{i=1}^{\myj-1} \myP^{(i)}}{\myj (\myj+1)}
\right).
\label{eq:8}
\end{equation}

with $\alpha = \frac{\myL \myMtot \myR_{\myc}}{\myK}$. 
Similarly, an upper bound $\bar{\myJ}_{\myUBrm}$ on the
average number of transmissions $\bar{\myJ}$ is given by
\begin{equation*}
    \bar{\myJ}_{\myUBrm}
    = \myE {\, \myJ \,\big|_{\myP_{\mypacrm}^{(\myj)} \leftarrow \myP^{(\myj)}} }
    = 1 + \sum_{\myj=1}^{\myJ_{\max}} \prod_{\myi=1}^{\myj} \myP^{(\myi)} .
\end{equation*}

\end{minipage}

\clearpage

\subsection{Cross-Layer Optimization Problem}
\label{sec:3c}
\myCLO\ is proposed to maximize $\myS_{\myLBrm}$ in (\ref{eq:8}) by carefully selecting
$(\mypbf, \mymbf, \myR_{\myc}) \in \Omega_{\myp} \times \Omega_{\mym} \times \Omega_{\myR_{\myc}}$
taking into account specific \myQoS\ constraints outlined below. Here, $\Omega_{\myp}$,
$\Omega_{\mym}$ and $\Omega_{\myR_{\myc}}$ represent the set $\myRbb^{+}_{0}$ of powers,
the set of $\myK$-length vectors of possible BL patterns \cite{Mansour2018_alt} and the set of
rate-matched \myLDPC\ code rates available at the \mySU\ \myTX{}. The maximization problem reads
\begin{equation*}
\left[ \hat{\mypbf}_{\myoptrm},\,
       \hat{\mymbf}_{\myoptrm},\,
       \hat{\myR}_{\myc,\myoptrm}
\right]
=\arg\hspace{-4ex}\max\limits_{(\mypbf,\,\mymbf,\,\myR_\myc)\in\Omega_{\mypbf}\times\Omega_{\mymbf}\times\Omega_{\myR_\myc}}
\hspace{-0.5ex}S_{\myLBrm}(\mypbf,\,\mymbf,\,\myR_\myc)
\end{equation*}

\begin{align}\hspace{2.5cm}
    \text{subject to}\qquad\quad 
    \myp_\myk &\ge 0                                   
        && \label{eq:9} \\
    \sum_{\myk\in\Omega_{0,\myk}} \myp_\myk &\le \myp_{\mytotrm}      
        && \label{eq:10} \\
    (\mym_\myk - 2)(\mym_\myk - 4)(\mym_\myk - 6) &= 0            
        && \label{eq:11} \\
    \sum_{\mym_\myk\in\Omega_{0,\myk}} \mym_\myk &= \myMtot      
        && \label{eq:12} \\
    \Pr\{\myI_\myk \le \myi_{\mytemprm}\} &\ge \mya
        && \label{eq:13} \\
    \myD &\le \myD_{\max}                             
        && \label{eq:14} \\
    L\myMtot \myR_\myc &\ge \myM_{\myreqrm} 
        && \label{eq:15} \\
    \prod_{\myj=1}^{\myJ_{\mymax}+1} \myP^{(\myj)} &\le \myp_{\myreqrm}.
        && \label{eq:16}
\end{align}
In~(\ref{eq:9}), the \myCLO{} allocates a non-negative $\myp_\myk$ to each subcarrier,
while adhering to a constraint on the total available \mySU{} \myTX{}
transmit power $\myp_{\mytotrm} \ge 0$ specified in~(\ref{eq:10}). The \myBL\ using the available \myQAM\ orders mentioned in Sect.\myII\ is represented in (\ref{eq:11}) maintaining the bit-rate constraint in~\!(\ref{eq:12}) with \!$\myMtot\!=2\myM$ and a fixed integer value $\myM$ chosen to support the intended
bit-rate with $\myK \le \myM \le 3K$ according to the available bit
patterns in~(\ref{eq:11}). In~(\ref{eq:13}), we ensure that the \mySU{} transmission
interference at the \myPU{} \myRX{} does not exceed a predefined $\myi_{\mytemprm}$
with a confidence level $\mya$. The closed form of the interference
constraint in~(\ref{eq:13}) is calculated similar to~\cite{Mansour2018_alt} according to
\begin{equation*}
\Pr\{\myI_\myk \le \myi_{\mytemprm}\}
    = 1 - \sum_{\myi=1}^{\myK_{\mym_\myk}}
        \psi\!\left(\myS_{\myi,\mym_\myk},\, \myd^{2}_{\myi,\mym_\myk},\, \myp_\myk,\, \myk\right) ,
\end{equation*}
where $\myI_\myk$ is the interference of the non-\myPU\ signal at the \myPU\ \myRX\ at the $\myk$th subcarrier. Furthermore,
{\small
\begin{equation*}
\psi\!\left(\myS, \myd^2, \myp, \myk\right)
= \myS \sqrt{\frac{4\phi_{\myp}}{\myp \sigma_{\mySrm\myPrm}^{2} d^{2}
    + \sigma_{\mathrm{\myP,\myI},\myk}^{2}}}\;
\myK_{1}\!\left(
        \sqrt{\frac{4\phi_{\myp}}{
        \myp \sigma_{\mySrm\myPrm}^{2} \myd^{2}
        + \sigma_{\myPrm,\myImath,\myk}^{2}}}
     \right),
\end{equation*}
}
with
$\phi_{\myp} = \frac{\myi_{\mytemprm}}{\myp\,\sigma_{\myPrm\myPrm}^{2}}$
and $\myK_{1}(\cdot)$ being the first-order modified Bessel function of the second kind.
Here,
$\sigma_{\myPrm\myPrm}^{2}
    = \myE{ \left| \myH_{\myPrm\myPrm,\myk,\myj} \right|^{2} }
\quad\text{and}\quad
\sigma_{\mySrm\myPrm}^{2}
    = \myE{ \left| \myH_{\mySrm\myPrm,\myk,\myj} \right|^{2}}$ are the channel gains between the \myPU\ \myTX\ and \myPU\ \myRX\, and the
\mySU\ \myTX\ and \myPU\ \myRX\, respectively. 
Furthermore, the \myCLO{} ensures that the \mySU{}
transmission satisfies the \myQoS\ constraints in (\ref{eq:14})-(\ref{eq:16}). First, by selecting
$\myJ_{\max} = \left\lfloor \myD_{\max}/\myD_{\mytriprm}\right\rfloor$ where $\myD_{\mytriprm}$ defines the round-trip delay of each
retransmission and $\myD_{\max}$ is the maximum delay, we ensure the total
transmission delay $\myD$ does not exceed $\myD_{\max}$ in (\ref{eq:14}).
Secondly, a suitable \linebreak
\newpage
$\myR_\myc \in \Omega_{\myR_\myc}$ is selected to guarantee a minimum required data
rate $\myM_{\myreqrm}$ in (\ref{eq:15}). Finally, in (\ref{eq:16}), the \myCLO\ ensures that the
\mySU\ \myPEP\ even after $\myJ_{\max}+1$ transmissions does not exceed a predefined threshold $\myp_{\myreqrm}$.



